\section*{الفصل \ref{ch:b2:heart} --- القلب والدورة القلبية}

\begin{solution}{exo:b2:heart:1}
الأجوفان $\to$ \omterm{def:b2:heart:anatomy}{الأذين} الأيمن $\to$ الدسام المثلث الشرف $\to$ \omterm{def:b2:heart:anatomy}{البطين}
الأيمن $\to$ الدسام الرئوي $\to$ \omterm{def:b2:blood-circulation:vessels}{الشريان} الرئوي $\to$ الرئتان
$\to$ الأوردة الرئوية $\to$ \omterm{def:b2:heart:anatomy}{الأذين} الأيسر $\to$ الدسام التاجي $\to$ \omterm{def:b2:heart:anatomy}{البطين}
الأيسر $\to$ الدسام الأبهري $\to$ الأبهر.
\end{solution}

\begin{solution}{exo:b2:heart:2}
الامتلاء (الانبساط): الدسامان الأذينيان البطينيان مفتوحان، والهلاليان
منغلقان. والانقباض متساوي الحجم: الأربعة كلها منغلقة. والقذف:
الهلاليان مفتوحان، والأذينيان البطينيان منغلقان. والارتخاء متساوي
الحجم: الكل منغلق. والصوت الأول انغلاق الدسامين الأذينيين
البطينيين في مستهل الانقباض؛ والثاني انغلاق الهلاليين في منتهاه.
\end{solution}

\begin{solution}{exo:b2:heart:3}
\omterm{prop:b2:heart:conduction}{العقدة الجيبية الأذينية} (عند $t = 0$؛ وهي الموجة P إذ يستقطب
الأذينان)؛ \omterm{prop:b2:heart:conduction}{والعقدة الأذينية البطينية} (تبلغها الدفعة عند \qty{80}{ms}
نحوًا، فتحبسها \qty{100}{ms}: وهي القطعة PR)؛ وحزمة هيس \omterm{prop:b2:heart:conduction}{وألياف
بركنجي} (\qtyrange{180}{220}{ms}؛ وهو المركّب QRS إذ يستقطب
البطينان)؛ ثم الموجة T بعد ذلك إذ يعيدان الاستقطاب.
\end{solution}

\begin{solution}{exo:b2:heart:4}
يقذف \omterm{def:b2:heart:anatomy}{البطين}، في مجال عمله، أكثر كلما امتلأ أكثر.
فإن ضخ \omterm{def:b2:heart:anatomy}{البطين} الأيمن أكثر، تلقى الأيسر أكثر بعد نبضات
قليلة فضخ أكثر، بالقانون نفسه: فيصحّح أي اختلال
نفسه بنفسه من غير إشارة.
\end{solution}

\begin{solution}{exo:b2:heart:5}
\omterm{prop:b2:heart:cycle}{حجم الضربة} \qty{80}{mL}؛ \omterm{prop:b2:heart:cycle}{وكسر القذف} $80/130 = \qty{62}{\%}$؛
والنتاج $80\times 70 = \qty{5.6}{L/min}$. وفي القصور: \qty{60}{mL}،
و$60/180 = \qty{33}{\%}$، و\qty{4.2}{L/min} --- \omterm{def:b2:heart:anatomy}{فالبطين} المتوسع
يبقي النتاج قريبًا من مستواه بعمله عند حجم كبير، لكنه
لا يقذف إلا ثلث ما يحمل.
\end{solution}

\begin{solution}{exo:b2:heart:6}
$W = 100\times 133\times 7\times 10^{-5} = \qty{0.93}{J}$؛ والاستطاعة
$0.93\times 1.2 = \qty{1.1}{W}$. وعند \qty{150}{mmHg}: \qty{1.4}{J}،
و\qty{1.7}{W}. والزيادة \qty{0.56}{W} ميكانيكية تساوي \qty{3.7}{W}
استقلابية، أي \qty{0.19}{mL} من الأكسجين في الثانية: \qty{11}{mL/min} زيادةً.
\end{solution}

\begin{solution}{exo:b2:heart:7}
$60/0.5 = 120$ نبضة في الدقيقة؛ و$60/1.5 = 40$. والثالث حصار
قلبي تام --- فالأذينان ينبضان عند 100 تحت العقدة، والبطينان عند 33
تحت ناظم خطى خاص بهما، \omterm{prop:b2:heart:conduction}{والعقدة الأذينية البطينية} لا توصّل شيئًا.
\end{solution}

\begin{solution}{exo:b2:heart:8}
الانجراف $20/25 = \qty{0.8}{s}$، وبإضافة \qty{0.15}{s}: \qty{0.95}{s}، أي 63
نبضة في الدقيقة. والأستيل كولين: $25/18 = \qty{1.39}{s}$، وبإضافة $0.15$:
\qty{1.54}{s}، أي 39 في الدقيقة. والنورأدرينالين: $20/40 = \qty{0.5}{s}$،
وبإضافة $0.15$: \qty{0.65}{s}، أي 92 في الدقيقة.
\end{solution}

\begin{solution}{exo:b2:heart:9}
يتيح التأخر للأذينين أن يفرغا في البطينين قبل أن
ينقبض البطينان؛ ولولاه لانقبض الأذينان والبطينان معًا،
ولانغلق الدسامان الأذينيان البطينيان على بطينين نصف
ممتلئين، ولهبط \omterm{prop:b2:heart:cycle}{حجم الضربة} بمقدار إسهام الأذينين.
وإطالة الفترة PR لا تفعل غير تأجيل الانقباض البطيني
ولا تكلف شيئًا حتى يطول التأخر فتسقط نبضات أو
يخفق التوصيل بالكلية.
\end{solution}

\begin{solution}{exo:b2:heart:10}
في الراحة: $5000/45 = \qty{111}{mL}$؛ وفي الجهد: $\num{30000}/180 =
\qty{167}{mL}$. فقانون ستارلنغ يحوّل العود الوريدي الأكبر
في الجهد إلى حجم ضربة أكبر؛ وتضيف الأعصاب الودّية المعدل
والقابلية للانقباض (أي حجمًا أصغر عند نهاية الانقباض). وفوق 200
نحوًا يقصر الانبساط عن ملء \omterm{def:b2:heart:anatomy}{البطين} فيهبط \omterm{prop:b2:heart:cycle}{حجم
الضربة} أسرع مما يرتفع المعدل.
\end{solution}

\begin{solution}{exo:b2:heart:11}
$v = Q/A$: أي $250/3 = \qty{83}{cm/s}$ و$250/1 = \qty{250}{cm/s}$.
و$\Delta P = \tfrac12\times 1060\times v^{2}$: \qty{365}{Pa}
(\qty{2.7}{mmHg}) و\qty{3300}{Pa} (\qty{25}{mmHg}). ويولّد \omterm{def:b2:heart:anatomy}{البطين}
الضغط الزائد بتثخين جداره؛ \omterm{def:b2:heart:anatomy}{والبطين} الثخين الجاسئ يتجاوز
في النهاية ما تحتمله تروية شرايينه الإكليلية ويمتلئ
امتلاءً رديئًا، ثم يخفق.
\end{solution}

\begin{solution}{exo:b2:heart:12}
التلقائية تعطي النبضة من غير أي دخل؛ وقانون ستارلنغ
يطابق النتاج على العود ويوازن بين البطينين من غير
أي دخل؛ والأعصاب والأدرينالين لا تضبط غير المعدل
والقابلية للانقباض حول ذلك اللب المنظِّم لنفسه. \omterm{def:b2:heart:anatomy}{والقلب} المزروع
ينبض (عند 100 نحوًا، بلا كابح مبهمي)، ويلائم حجم ضربته
للامتلاء، ويستطيع رفع نتاجه في الجهد بآلية
ستارلنغ وبالأدرينالين الدائر --- لكن ببطء،
وأقل من \omterm{def:b2:heart:anatomy}{قلب} معصّب.
\end{solution}

\begin{solution}{pb:b2:heart:1}
\textbf{1.} \qty{70}{mL}؛ و$70/130 = \qty{54}{\%}$؛ و$70\times 70 =
\qty{4.9}{L/min}$.
\textbf{2.} $60/70 = \qty{0.86}{s}$؛ والانبساط $0.86 - 0.3 =
\qty{0.56}{s}$.
\textbf{3.} $70/0.3 = \qty{233}{mL/s}$؛ و$233/3 = \qty{78}{cm/s}$.
\textbf{4.} $70/1500 = \qty{0.047}{s}$، أي \qty{50}{ms} نحوًا.
\textbf{5.} $100\times 133\times 7\times 10^{-5} = \qty{0.93}{J}$؛
و$0.93\times 70/60 = \qty{1.1}{W}$.
\textbf{6.} البطينان معًا $1.1\times 1.2 = \qty{1.3}{W}$؛ وعند
\qty{15}{\%}: \qty{8.7}{W} استقلابية؛ و$8.7/20 = \qty{0.43}{mL}$ من
الأكسجين في الثانية، أي \qty{26}{mL/min}.
\textbf{7.} باستخلاص \qty{0.14}{mL} من كل مليلتر من الدم:
$26/0.14 = \qty{190}{mL/min}$.
\textbf{8.} الفترة \qty{0.86}{s}، وزمن الانجراف \qty{0.71}{s}:
$20/0.71 = \qty{28}{mV/s}$.
\textbf{9.} الفترة \qty{0.6}{s}، والانجراف \qty{0.45}{s}: \qty{44}{mV/s}.
\textbf{10.} الفترة \qty{0.33}{s}، والانجراف \qty{0.18}{s}: \qty{110}{mV/s}.
\textbf{11.} الفترة R--R \qty{0.86}{s}؛ وP استقطاب الأذينين؛ وQRS
استقطاب البطينين؛ وT إعادة استقطاب البطينين.
\textbf{12.} يتيح التأخر للأذينين أن يفرغا في البطينين
قبل أن ينقبضا؛ وعند 180 تصير الدورة كلها \qty{0.33}{s}، فتأخر
لا يتغير مقداره \qty{0.16}{s} يلتهم نصفها --- فالأعصاب
الودّية تسرّع توصيل العقدة أيضًا.
\textbf{13.} حصار قلبي تام: فالبطينان عند 40 في الدقيقة،
توقّتهما الحزمة أو \omterm{prop:b2:heart:conduction}{ألياف بركنجي}، والأذينان عند 75 تحت
\omterm{prop:b2:heart:conduction}{العقدة الجيبية الأذينية}.
\textbf{14.} \qty{120}{mL}؛ و$120/150 = \qty{80}{\%}$؛ و$120\times 180 =
\qty{21.6}{L/min}$.
\textbf{15.} $0.33 - 0.3 = \qty{0.03}{s}$: أي لا وقت يكاد يكفي للامتلاء؛
وأسرع من ذلك ينهار \omterm{prop:b2:heart:cycle}{حجم الضربة}.
\textbf{16.} $W = 120\times 133\times 1.2\times 10^{-4} = \qty{1.9}{J}$،
و$\times 3$ في الثانية: \qty{5.7}{W}؛ والبطينان معًا \qty{6.9}{W}؛
والاستقلابي \qty{46}{W}؛ والأكسجين \qty{2.3}{mL/s}، أي \qty{140}{mL/min}.
\textbf{17.} $140/0.14 = \qty{1000}{mL/min}$، أي خمسة أمثال جريان
الراحة، تُسلَّم في انبساط انكمش إلى عُشر الدورة:
فلا بد أن تتوسع \omterm{def:b2:blood-circulation:vessels}{الشرينات} الإكليلية إلى حدها.
\textbf{18.} المعدل 100 نحوًا بلا كابح مبهمي، والحجم عند نهاية
الانقباض لم يتغير عند \qty{60}{mL}: \omterm{prop:b2:heart:cycle}{فحجم الضربة} $150 - 60 = \qty{90}{mL}$،
والنتاج \qty{9}{L/min}.
\textbf{19.} المعدل وحده ($70 \to 180$ عند \qty{70}{mL}، أي
\qty{12.6}{L/min}): $+7.7$؛ والامتلاء (\qty{20}{mL} زيادةً عند 180):
$+3.6$؛ والقابلية للانقباض (\qty{30}{mL} أقل من المتبقي عند 180): $+5.4$ ---
من 4.9 إلى \qty{21.6}{L/min}، والإسهامات الثلاثة تبلغ
$+16.7$ بالضبط.
\textbf{20.} $233/1 = \qty{233}{cm/s}$، أي \qty{2.3}{m/s}.
\textbf{21.} $\tfrac12\times 1060\times 2.33^{2} = \qty{2900}{Pa}$، أي
\qty{22}{mmHg}.
\textbf{22.} $120/0.3 = \qty{400}{mL/s}$، أي \qty{4}{m/s}؛ و$\tfrac12\times
1060\times 16 = \qty{8500}{Pa}$، أي \qty{64}{mmHg}.
\textbf{23.} الذروة $120 + 22 = \qty{142}{mmHg}$ في الراحة، ونحو $140 +
64 = \qty{204}{mmHg}$ في الجهد؛ وضغط القذف المتوسط في
الراحة يرتفع من 100 إلى 122: أي شغل في النبضة $\times 1.22$.
\textbf{24.} الجدار الأثخن يولّد قوة أكبر ويخفض
الإجهاد على الليف الواحد؛ لكن الكتلة الواجب ترويتها تنمو بينما
الجريان الإكليلي، إذ يعصره الجدار الثخين في الانقباض ويُعطى
انبساطًا أقصر، لا يجاريها، والجدار الثخين يرتخي
رداءةً ويمتلئ أقل: فيتبع ذلك الإقفار والقصور.
\textbf{25.} \qty{4.9}{L/min} في الراحة، و\qty{21.6}{L/min} في الجهد؛
و\qty{0.93}{J} في النبضة؛ والتدرج \qty{22}{mmHg} في الراحة
و\qty{64}{mmHg} في الجهد.
\end{solution}
